% titel: Blatt 3, A7: Hess'scher Satz (Standard-Reaktionsenthalpie)
% skript: Mortimer Kap. 5.4 (Hess'scher Satz, Reaktionsenthalpie)

% Teil I: Werkzeuge, kein Lösungsweg.

Die Reaktionsenthalpie einer Reaktion erhält man, ohne einzelne Bindungen
aufzubrechen zu müssen: Man kombiniert bekannte Teilreaktionen so, dass sie
zur Zielreaktion aufsummiert werden, und addiert die Enthalpien.

\begin{werkzeugkasten}[Hess'scher Satz]
Die Enthalpieänderung einer Gesamtreaktion ist unabhängig vom Weg; sie ist die
Summe der Enthalpien der Teilreaktionen, mit denen man die Zielreaktion
zusammensetzt:
\begin{equation}
  \Delta_rH^\circ = \sum_i \nu_i\, \Delta H_i^\circ.
\end{equation}
Regeln beim Umformen von Teilreaktionen:
\begin{itemize}
  \item Eine Teilreaktion \emph{umdrehen} skaliert $\Delta H$ mit $-1$.
  \item Eine Teilreaktion mit einem Faktor multiplizieren skaliert $\Delta H$
        um denselben Faktor.
  \item Zwischenprodukte, die insgesamt auf beiden Seiten stehen, löschen sich
        heraus. Dazu Rektionen so wählen, dass die Edukte der Zielreaktion
        links und die Produkte rechts stehen bleiben.
\end{itemize}
\textit{Skript:} Mortimer Kap. 5.4 (Hess'scher Satz), Übung 3/7.
\end{werkzeugkasten}

\begin{achtung}[Darauf achten]
\begin{itemize}
  \item Entstehen Zwischenprodukte mit stöchiometrischen Faktoren, muss man
        die Teilreaktion entsprechend \emph{vervielfachen} --- und die
        Enthalpie mitziehen.
  \item Vorzeichen: beim Umkehren einer Teilreaktion dreht sich das Vorzeichen
        von $\Delta H$ um.
  \item Beim Streichen von Zwischenprodukten exakt spüren, dass Atom- und
        Massenbilanz zusammenpassen; am Ende muss die Zielreaktion stehen.
  \item Einheit aller Beiträge: kJ/mol (auf die Zielreaktion bezogen).
\end{itemize}
\end{achtung}
